\chapter{Discussion}
\section{The refractive index of air}
The biggest problem in the experiment was how we measured the pressures.
Firstly the pressure gauge inside the bicycle pump had $0.1$ bar increments, which is quite a lot, so even under ideal circumstances the measurement would not be great.
What made it worse was that we 'measured the pressure', by starting the measurement, while it was dropping, because a constant leak was the best way we found to maintain a somewhat managable flow.
Another huge problem with this method was that the flow slowed down as the pressure decreased, so the last air flowed out really slowly, 
so each measurement actually stopped at some point slightly above atmospheric pressure, that we could not measure because of the poor precision pressure gauge.\\
Considering the aforementioned problem of the pressure at $p_{0}$ and at $p$ being very uncertain we were still able to reach a refractive index of $n=1.000214 \pm 5.4 \cdot10^{-6}$.
This was still out of the interval of confidence. But if $p_{0}$ was increased by the $0.111 \text{bar}$ which is close to the $0.1 \text{bar}$ of uncertaincy we expected, we were able to get the expected of $n=1.00027717$ \autocite{air} within the interval of confidence.
\section{The validity of the intensity hypothesis}
As can be seen in figure \ref{fig:WaveplateFit}, 
we have achieved very great agreement between measurements and theory.
The residuals seem to follow a kind of trend, where positive and negative deviations are grouped together.
If we run a Shapiro-Wilks test on the residuals, we get a p value of 0.053, for the hypothesis of them following a normal distribution,
which would be within a significance level of 5\%.
We have reason to believe this is not the case however, 
as we observed during the experiment that the background intensity would vary, based on if someone was walking by, 
or someone changed the way they sit. 
Therefore it would make sense if deviations are grouped, as major changes might not happen after every measurement.
To combat this, the experiment should of course be done in a dark room, but that was not possible in this case.